% !TeX root = main.tex
\subsection{高温胁迫下黑麦草枯叶表型的全基因组关联分析}
\label{sec:ryegrass-heat-gwas}

\subsubsection{高温胁迫试验与表型获取}
\label{sec:ryegrass-heat-phenotyping}
试验在武汉植物园江夏安山基地开展。供试材料为课题组前期从全球范围收集并完成基因组重测序的 223 份多年生黑麦草（\textit{Lolium perenne} L.）种质。每份种质由单粒种子萌发获得初始植株，再经无性繁殖扩繁，以保证用于田间表型测定的植株与对应的重测序材料具有相同的遗传背景。供试材料于 2025 年 12 月 30 日种植，每份种质设置 3 个生物学重复。\par
2026 年 6 月 10 日，逐株调查叶片枯黄程度并采集植株图像，分别获取人工枯叶率和图像枯叶率。人工枯叶率由调查人员目测估计，以百分数记录。图像枯叶率按第~\ref{sec:ryegrass-senescence}~节所述方法计算：先利用 BiMSeg 分割植株与背景，再由 XGBoost 对植株掩膜内的像素进行绿叶与枯叶分类，以枯叶像素数占植株前景像素总数的百分比表示可见冠层的枯叶率。对每份种质，分别计算 3 个生物学重复的人工枯叶率均值和图像枯叶率均值，作为该种质的两项表型观测值。\par
按种质编号将两项表型记录与重测序 VCF 文件中的样本进行匹配，最终有 197 份种质同时具备人工枯叶率、图像枯叶率及基因型数据。以这 197 份种质为共同样本，对两种枯叶率进行配对比较，并分别以人工枯叶率和图像枯叶率开展全基因组关联分析。

\subsubsection{枯叶率的分布与测量一致性}
\label{sec:ryegrass-heat-phenotype-comparison}
以同时具有表型与基因型记录的 197 份多年生黑麦草种质为对象，比较高温胁迫下图像枯叶率与人工枯叶率的群体分布（图~\ref{fig:ryegrass-heat-phenotype-distributions}），并按种质编号对两种测量值进行配对分析。
\begin{figure}[H]
  \centering
  \includegraphics[width=0.95\linewidth]{figures/heat_phenotype_distributions.pdf}
  \methodbicaption{高温胁迫下多年生黑麦草图像枯叶率与人工枯叶率的分布}
    {Distributions of image-derived and visually assessed senescent leaf percentages in perennial ryegrass under heat stress}
    {fig:ryegrass-heat-phenotype-distributions}
  \par\smallskip
  {\raggedright\normalfont\zihao{5}
   注：a、b 分别为图像枯叶率与人工枯叶率的直方图，虚线表示拟合的正态分布曲线；c、d 为相应的正态 Q--Q 图。\par
   Note: Panels a and b show histograms of image-derived and visually assessed senescent leaf percentages, respectively. Dashed curves indicate fitted normal distributions. Panels c and d show the corresponding normal Q--Q plots.\par}
\end{figure}
图像枯叶率和人工枯叶率的群体均值分别为 72.30\% 和 72.60\%，但两种测量值的分布有所不同。图像枯叶率介于 34.13\% 至 100.00\%，中位数为 73.03\%；人工枯叶率介于 10.00\% 至 100.00\%，中位数为 80.00\%。与图像测量相比，人工测量的中位数较高，分布却向低值端延伸得更远，因而其均值明显低于中位数。两种表型均呈不同程度的左偏，人工枯叶率的左偏更为明显。Shapiro--Wilk 检验表明，两种表型均偏离正态分布（图像枯叶率：\(W=0.954\)；人工枯叶率：\(W=0.914\)；均为 \(P<0.001\)）。\par
群体均值反映总体枯黄水平，但不能说明同一份种质的两种测量值相差多少。为此，进一步采用配对散点图比较两种测量值的变化趋势，并通过 Bland--Altman 差值图考察种质水平上的测量一致性（图~\ref{fig:ryegrass-heat-phenotype-agreement}）。
\begin{figure}[H]
  \centering
  \includegraphics[width=0.95\linewidth]{figures/heat_phenotype_agreement.pdf}
  \methodbicaption{高温胁迫下多年生黑麦草图像枯叶率与人工枯叶率的配对比较}
    {Paired comparison of image-derived and visually assessed senescent leaf percentages in perennial ryegrass under heat stress}
    {fig:ryegrass-heat-phenotype-agreement}
  \par\smallskip
  {\raggedright\normalfont\zihao{5}
   注：每个点代表一份种质。a 为配对散点图，虚线表示两种测量值相等；b 为 Bland--Altman 差值图，横坐标为两种测量值的均值，纵坐标为图像枯叶率减去人工枯叶率的差值。实线表示平均差，虚线表示估计的 95\% 一致性界限。\par
   Note: Each point represents one accession. In panel a, the dashed line indicates equality between the two measurements. Panel b plots the paired difference (image-derived minus visually assessed) against the mean of the two measurements. The solid line indicates the mean difference, and the dashed lines indicate the estimated 95\% limits of agreement.\par}
\end{figure}
图像枯叶率与人工枯叶率呈正相关（Pearson \(r=0.813\)，Spearman \(\rho=0.776\)）。人工目测枯黄程度较重的种质，其图像枯叶率通常也较高，表明图像测量能够反映种质间枯黄程度的相对差异。以图像枯叶率减去人工枯叶率，平均差为 \(-0.30\) 个百分点，而配对差值绝对值的均值为 10.12 个百分点，估计的 95\% 一致性界限为 \(-26.16\) 至 \(25.56\) 个百分点。可见，两种测量对群体平均枯黄水平的估计接近，但逐份种质比较时，正、负差值在计算平均差时相互抵消，单份种质的测量差异仍较明显。





